\section{El cuello de botella de los datos en robótica: el mundo real es demasiado complejo y los datos, demasiado caros}

Ya se dijo que el modelo todavía no es completamente independiente; esta sección entra en el problema más duro de la robótica: los datos. Tan Jie (谭杰) considera que el mayor problema de los robots son los datos. Los robots operan en un unstructured environment complejo, donde puede pasar cualquier cosa; las trayectorias de acción son caras, los fallos de cola larga son difíciles de cubrir y, además, las diferencias entre cuerpos de hardware dificultan compartir los datos.

\begin{figure}[H]
\centering
\includegraphics[width=0.96\textwidth]{robotics-data-bottleneck.png}
\caption{El cuello de botella de los datos en robótica: los entornos no estructurados, los fallos de cola larga y la retroalimentación de las acciones encarecen especialmente los datos. Diagrama conceptual propio, elaborado a partir de la conversación entre 00:23:44 y 00:27:52.}
\end{figure}

\begin{knowledgebox}{Lectura de la figura: los datos de robótica no son datos de video comunes}
El entorno real no es controlable, la captura de trayectorias de acción es cara, las muestras de fallo son difíciles de cubrir, entre cuerpos distintos hay diferencias de hardware y la tasa de éxito en la manipulación fina es baja. Por eso los datos de alta calidad se convierten en el cuello de botella central para que los robots pasen de la demo a la generalización.
\end{knowledgebox}

\subsection{Por qué unos cientos de horas no bastan para verificar el scaling}

Esta subsección explica por qué el problema de los datos en robótica es más difícil que el de los datos comunes de video o de texto. Entrenar un robot no consiste solo en reunir videos que «parecen relacionados», sino en reunir datos que incluyan acciones, resultados, fallos y rutas de recuperación.

Tan Jie menciona que, si en robótica solo se tienen unos cientos de horas de datos, quizá no sea posible verificar el scaling; podrían necesitarse decenas de miles de horas de datos para ver si la recipe de entrenamiento realmente funciona. Esta exigencia es muy dura tanto para los equipos de emprendimiento como para los inversionistas, porque al principio, con poco dinero y pocos datos, es muy difícil demostrar una ruta de largo plazo, y la robótica necesita justamente una convicción de largo plazo y una gran inversión.

\begin{warningbox}{Lo engañoso de las demos con pocos datos}
Con decenas o cientos de horas de datos es posible lograr una demo vistosa, pero eso no basta para demostrar generalización. Lo que la robótica necesita en realidad son datos de alta calidad que abarquen distintos escenarios, objetos, cuerpos y modos de fallo.
\end{warningbox}

\subsection{Por qué la ruta de Tesla es difícil de replicar en los robots de propósito general}

Los automóviles de Tesla son útiles, así que todos los días generan datos reales de conducción y el ciclo virtuoso de datos puede ponerse en marcha. Muchos equipos de hardware robótico pueden capturar datos, pero todavía no alcanzan el umbral de «ser útiles por sí mismos», por lo que su ciclo virtuoso de datos no arranca. Por eso no se puede afirmar sin más que «China tiene hardware y, por lo tanto, puede replicar la ruta de Tesla»: primero el hardware debe entrar en uso real para que los datos regresen de forma continua.

\subsection{Los robots son como niños pequeños: capacidades desiguales}

Esta subsección añade un juicio que se pasa por alto con facilidad: las capacidades de los robots no crecen de manera uniforme. Tan Jie recurre a la analogía con un niño: considera que la locomotion de los robots ya es muy sólida, e incluso algunas capacidades motrices superan a las de un adulto; pero la manipulation, sobre todo la manipulación con manos diestras, quizá siga al nivel de un niño de dos o tres años: entiende las instrucciones de forma aproximada, lo intenta varias veces, pero no sujeta los objetos con firmeza, y los movimientos finos le cuestan todavía más.

\begin{knowledgebox}{Capacidades desiguales: Locomotion vs Manipulation}
\begin{tabularx}{\textwidth}{p{0.24\textwidth}p{0.34\textwidth}X}
\toprule
Capacidad & Estado actual & Por qué no avanzan al mismo ritmo \\
\midrule
Locomotion & Correr, saltar, el equilibrio y la marcha ya han avanzado muy rápido & En los últimos cinco años, el aprendizaje por refuerzo y Sim2Real resolvieron en lo esencial muchos problemas de control motor. \\
Sujeción común & Puede tomar y colocar objetos sencillos y seguir instrucciones & Requiere visión, comprensión del objetivo y control grueso, pero la tolerancia a errores es relativamente alta. \\
Manipulación con manos diestras & Sigue siendo muy difícil; muchas tareas son inestables & La coordinación de varios dedos, el tacto, la fuerza de contacto y la retroalimentación del material son muy complejos. \\
Planificación cognitiva & Los modelos grandes aportan sentido común y descomposición de tareas & Pero el plan sigue dependiendo de la tasa de éxito de las acciones reales. \\
\bottomrule
\end{tabularx}
\end{knowledgebox}

\subsection{Resumen de la sección}

El problema de primeros principios de la robótica sigue siendo contar con datos de alta calidad. Sin datos suficientemente reales, diversos y verificables, no es posible juzgar si el modelo puede escalar ni pasar de la demo a un producto utilizable.
